\chapter{Модель данных и рабочего процесса}

\section{Кадр, серия и объект}

Начиная с 0.5.52.0 модель явно разделяет физическую экспозицию и её файловое представление. \textbf{Observation} --- одна физическая экспозиция с path-independent identity; \textbf{кадр (frame)} --- конкретный FITS/file product, поэтому несколько frames могут ссылаться на один observation. \textbf{Серия} --- набор кадров, относящихся к одному согласованному наблюдательному контексту. \textbf{Объект} --- устойчивый источник, идентифицируемый через каталог и БД.

Legacy \texttt{frames.frame\_uid} сохраняется и не переписывается. Новая связь \texttt{frames.observation\_id} вводится отдельно, чтобы копии или processing-варианты одной экспозиции не требовали объединения существующих frame rows. Разделение этих уровней принципиально: file identity, physical exposure identity, калибровка одной серии и долговременная межсерийная калибровка объекта являются разными понятиями.

\section{Calibration products}

Современная архитектура сохраняет базовые measurements и нормализованные calibration terms. Промежуточные калибровки имеют provenance и immutable generations. Финальный пользовательский слой --- science measurements, восстановленные из базовой фотометрии и опубликованных calibration terms.

\section{Почему не следует просто переписывать measurements}

Хранение calibration generations позволяет:

\begin{itemize}
  \item повторить обработку новой версией алгоритма;
  \item сравнить две calibration generations;
  \item не уничтожать предыдущий результат;
  \item восстановить, каким именно набором terms была получена science magnitude.
\end{itemize}

\section{Manifest и pipeline run}

Каждый recorded pipeline run хранит structured steps и их результаты. Manifest предназначен для пользовательского и машинного контроля запуска; таблицы pipeline runs позволяют сравнивать этапы и производительность между версиями.
